\documentclass[12pt]{article}
\usepackage[utf8]{inputenc}
\usepackage[T1]{fontenc}
\usepackage{hyperref}
\title{Zhaba Language Resources for K-12 Curriculum Development}
\author{Compiled by Claude Sonnet 5.5 (claude-sonnet-5-5) for Andr\'as Kornai}
\date{}
\begin{document}
\maketitle

\section{Introduction}
Zhaba (ISO 639-3 zhb; Glottocode zhab1238; also called nDrapa) is a Qiangic Sino-Tibetan language of Sichuan, China. It has about 7,800 speakers. Speakers live mainly in the Xianshui River valley, in Daofu and Yajiang counties. Searches found only a small body of scholarly description. They found no school materials, primers or dictionaries. This report lists only what the searches returned.

\section{Primary Educational Resources}
No primers, textbooks or school programmes in Zhaba turned up in the searches, and Zhaba is not known to be taught in schools. A curriculum would have to be built from scratch, drawing on the linguistic descriptions below and working with the community.

\section{Linguistic and Reference Materials}
\begin{itemize}
\item Wikipedia overview of the Zhaba language, with the phonology, dialects and speaker numbers: \url{https://en.wikipedia.org/wiki/Zhaba_language}
\item Glottolog reference record pointing to Gong Qunhu's \emph{Zhabayu Yanjiu} (A Study of the nDrapa Language, 2007), the main full-length description: \url{https://glottolog.org/resource/reference/id/cldf:gong:07:2}
\item Endangered Languages Project bibliography for Zhaba: \url{https://www.endangeredlanguages.com/lang/5610/bibliography}
\item ASJP word-list data for Zhaba: \url{https://asjp.clld.org/languages/ZHABA}
\end{itemize}

\section{Religious and Community Resources}
No scripture translations or community organisations for Zhaba were found.

\section{Digital and Audio Resources}
No audio, video or software resources specific to Zhaba were found. The ASJP and Endangered Languages Project pages above are the only online data points found.

\section{Curriculum Development Context}
Zhaba is used as a first language by all adults but not by all young people, so intergenerational transmission is weakening. A curriculum would need new documentation, such as a practical orthography, a lexicon and recorded texts, made with native speakers. The scholarly descriptions above could inform that work. Little else is available.

\end{document}